% META: {"id":"1SPE-PROBCOND-EX-004","chapitre":"1SPE-PROBA-COND","type_objet":"exercice","capacites":["1SPE-PROBA-COND-C1"],"capacites_codes":["C1"],"methodes":["M1"],"parcours":1,"competences":["calculer","raisonner"],"duree_min":10,"mode_creation":"ex_nihilo","sources_inspiration":[],"parametres_sympy":{},"fichier_tex":"chapitres/1SPE-PROBA-COND/exercices/1SPE-PROBCOND-EX-004.tex","corrige_tex":"chapitres/1SPE-PROBA-COND/corriges/1SPE-PROBCOND-CO-004.tex","status":"approved","origin":"generated","status_history":["generated","audited","accepted","approved"],"release_acceptance":"RELEASE_OWNER_BATCH_ACCEPTANCE","acceptance_closure_digest":"sha256:9b3ccf9a81c5520fbb7e03b7b2d3e7bf2057a02834b6d908bf3be5f49f3b553f","acceptance_receipt_digest":"sha256:4fad86f0282e0fa4e0419cca1ad6379ae7af5a280c04a4a90880f90a1c044242"}
% BEGIN-VERIFY
% from sympy import Rational
% P_A = Rational(1, 4)
% P_A_B = Rational(3, 5)
% P_AinterB = P_A * P_A_B
% assert P_AinterB == Rational(3, 20)
% END-VERIFY
\begin{exercice}{1SPE-PROBCOND-EX-004}{1}{10}
On donne $P(A) = \dfrac{1}{4}$ et $P_A(B) = \dfrac{3}{5}$.

\begin{enumerate}
  \item Calculer $P(A \cap B)$.
  \item Si de plus $P(B) = \dfrac{1}{2}$, calculer $P_B(A)$.
\end{enumerate}
\end{exercice}
